% Chapter 9, typeset from lectures/L09: the README is the opener and the closing review, the
% lecture's Appendix A is a section with the same subsections, and its Appendix B, the exercises,
% closes the chapter. The gpio_driver_demo is typeset straight from its files in §9.1.5.

\chapter{Structs as Drivers}
\label{c9:ch}

{\large\itshape One struct per instance, and one set of functions for all of them.\par}

\begin{chapterbox}{In this chapter}
\begin{itemize}
  \item Why Chapter~\ref{c4:ch}'s ADC and Chapters~\ref{c6:ch}--\ref{c8:ch}'s
    UART/EEPROM/watchdog functions got away with no struct (they're singletons), and why
    per-instance macros and globals stop scaling for everything else: revisiting
    Chapter~\ref{c4:ch}'s multi-LED exercise.
  \item Bundling a peripheral's state into a struct, and writing ``methods'' as functions that
    take a pointer to that struct as their first argument.
  \item Building a GPIO (\code{gpio_t}) driver, split into a header and a source file.
  \item Worked example: rebuilding Chapter~\ref{c8:ch}'s watchdog demo on \code{gpio_t}, adding
    pin-change-interrupt support to the GPIO driver along the way.
  \item Exercise: a complete \code{timer_t} driver, with its own header/source split and a
    callback, built the same way as \code{gpio_t}.
\end{itemize}
\end{chapterbox}

Every driver so far has served exactly one of something: one ADC, one UART, one EEPROM, one
watchdog, and LEDs and buttons hardcoded into macros one pin at a time. This chapter is about what
to do when a program needs several of the same thing. After it you should be able to explain why
hardcoded per-instance functions or macros break down once a program needs more than one of
something, design a struct that bundles a peripheral's configuration and state, write ``methods''
as plain functions taking a pointer to that struct as their first parameter (the \code{self}
convention), and refactor a program built on duplicated globals into one built on an array of
struct instances.

Read \secref{c9:sec:structs}, the substance of the chapter, in order: it builds a GPIO driver one
step at a time and ends with Chapter~\ref{c8:ch}'s demo rebuilt on it. Then work the exercises in
\secref{c9:sec:exercises}. The Reflection exercises, and the design questions the Code exercises
ask, are answered in Appendix~\ref{app:answers}; the Code exercises have no printed solution.

% Section 9.1, from lectures/L09/appendix/a_structs_as_drivers.md. The gpio_driver_demo's GPIO
% driver and main.c are typeset straight from their files in §9.1.5; its EEPROM, serial and
% watchdog drivers are the ones printed in Chapters 6 to 8.

\section{Structs as Drivers}
\label{c9:sec:structs}\label{c9:app:a}

\subsection{Why bundle state}
\label{c9:a:1}
\begin{itemize}
  \item The drivers covered in Chapters~\ref{c4:ch} and \ref{c6:ch}--\ref{c8:ch} --- the ADC,
    UART, EEPROM, and the watchdog --- did not need any structs:
  \begin{itemize}
    \item The chip has exactly one ADC (\code{adc_read()} already takes a channel number as a
      parameter, so it was never hardcoded to a single pin the way the examples below are).
    \item It has exactly one UART.
    \item It has exactly one EEPROM.
    \item It has exactly one watchdog.
    \item Because there is only one instance of each peripheral, plain functions operating on
      implicit global state were the simplest correct solution.
  \end{itemize}
  \item The same single-instance pattern shows up again with peripherals the chip actually has
    many of, where it isn't the right solution:
  \begin{itemize}
    \item \code{LED_ON} assumes one specific LED, even though the chip has many GPIO pins.
    \item The configuration is effectively hardcoded into macros, functions, and global state.
  \end{itemize}
  \item This pattern works well as long as the program needs only one instance of each thing.
  \item The pattern starts to break down when the program needs several instances:
  \begin{itemize}
    \item \Secref{c4:app:b} contains an exercise with five LEDs, Exercise~\ref{c4:ex:7}.
    \item Without a better abstraction, adding another LED means copying an entire family of
      macros or functions and renaming them.
    \item That quickly becomes repetitive, difficult to maintain, and easy to get wrong.
  \end{itemize}
  \item C already provides the tool needed to solve this problem: the \code{struct}:
  \begin{itemize}
    \item Store one peripheral instance's configuration and state in a \code{struct}.
    \item Create one struct value for each LED (or other multi-instance peripheral) instance.
    \item Write ordinary functions that receive a pointer to the relevant struct.
    \item The same functions can then operate on several different instances.
  \end{itemize}
  \item By convention, the pointer parameter is called \code{self}:
  \begin{itemize}
    \item \code{self} points to the particular peripheral instance that the function should
      operate on.
    \item It plays a role similar to \code{this} in languages with built-in object support, such
      as C++, C\#, and Java.
    \item Unlike \code{this}, however, \code{self} is just an ordinary function parameter written
      explicitly in the code.
  \end{itemize}
  \item However, C has no built-in support for object-oriented programming:
  \begin{itemize}
    \item There is no hidden behaviour.
    \item There are no constructors.
    \item There is no inheritance.
    \item There is no automatic method dispatch.
    \item There is only a \code{struct} containing data and a set of ordinary functions that
      accept a pointer to that data.
  \end{itemize}
  \item This chapter keeps the design deliberately simple:
  \begin{itemize}
    \item The struct's fields remain directly accessible.
    \item Functions are called normally.
    \item No function pointers are required.
  \end{itemize}
  \item The next chapter, Chapter~\ref{c10:ch}, explores how much further this approach can be
    taken:
  \begin{itemize}
    \item Hiding a struct's internal fields.
    \item Calling operations through function pointers.
    \item Creating a more object-like interface while still using C.
  \end{itemize}
\end{itemize}

\subsection{The \code{self} convention: a \code{GPIO} driver}
\label{c9:a:2}
A minimal driver is a struct holding whatever a function needs to know to act on one specific
instance: here, pointers to the GPIO's data direction (\code{DDRx}), output (\code{PORTx}), and
input (\code{PINx}) registers, plus the bit position of its pin within them.

See \secref{c1:a:3} for a refresher on structs and \code{typedef}; \code{gpio_t} follows the same
anonymous-struct-plus-\code{_t}-suffix convention introduced there.

\begin{ccode}
/**
 * @brief GPIO driver structure.
 */
typedef struct
{
    /** Pointer to the data direction register associated with the GPIO. */
    volatile uint8_t* ddrx;

    /** Pointer to the output register associated with the GPIO. */
    volatile uint8_t* portx;

    /** Pointer to the input register associated with the GPIO. */
    volatile uint8_t* pinx;

    /** Bit position of the GPIO pin within the associated port. */
    const uint8_t pin;
} gpio_t;

/**
 * @brief Enumeration of GPIO modes.
 */
typedef enum
{
    GPIO_MODE_INPUT,        ///< Standard GPIO input.
    GPIO_MODE_INPUT_PULLUP, ///< GPIO input with its internal pull-up enabled.
    GPIO_MODE_OUTPUT,       ///< GPIO output.
} gpio_mode_t;

/**
 * @brief Initialize GPIO.
 *
 * @param[in] self GPIO instance.
 * @param[in] mode GPIO mode.
 */
void gpio_init(gpio_t* self, const gpio_mode_t mode)
{
    if (NULL == self) { return; }

    switch (mode)
    {
        case GPIO_MODE_INPUT_PULLUP:
            *self->ddrx &= ~(1U << self->pin);
            *self->portx |= (1U << self->pin);
            break;
        case GPIO_MODE_OUTPUT:
            *self->ddrx |= (1U << self->pin);
            *self->portx &= ~(1U << self->pin);
            break;
        default:
            *self->ddrx &= ~(1U << self->pin);
            *self->portx &= ~(1U << self->pin);
            break;
    }
}

/**
 * @brief Set state of the given GPIO.
 *
 * @param[in] self GPIO instance.
 * @param[in] state GPIO state (true = high, false = low).
 *
 * @note This operation is only supported for GPIOs configured as outputs.
 */
void gpio_write(gpio_t* self, const bool state)
{
    if (NULL == self) { return; }
    const bool output = (*self->ddrx & (1U << self->pin));
    if (!output) { return; }

    if (state) { *self->portx |= (1U << self->pin); }
    else { *self->portx &= ~(1U << self->pin); }
}

/**
 * @brief Read state of the given GPIO.
 *
 * @param[in] self GPIO instance.
 *
 * @return GPIO state (true = high, false = low).
 */
bool gpio_read(const gpio_t* self)
{
    return (NULL != self) ? (*self->pinx & (1U << self->pin)) : false;
}

/**
 * @brief Toggle state of the given GPIO.
 *
 * @param[in] self GPIO instance.
 */
void gpio_toggle(gpio_t* self)
{
    if (NULL == self) { return; }
    const bool output = (*self->ddrx & (1U << self->pin));
    if (!output) { return; }
    *self->pinx = (1U << self->pin);
}
\end{ccode}

\noindent A GPIO output instance can then be initialized as shown below:

\begin{ccode}
#define LED1_PIN 1U // D9 -> PORTB1.

/** Initialize LED on D9 (PORTB1). */
gpio_t led1 = {&DDRB, &PORTB, &PINB, LED1_PIN};

gpio_init(&led1, GPIO_MODE_OUTPUT);
\end{ccode}

\noindent Most of these functions take \code{gpio_t*} rather than \code{gpio_t}: passing a pointer
avoids copying the struct on every call, and they need to write through \code{self} to the
hardware registers it points at. \code{gpio_read()} takes \code{const gpio_t*} instead, since it
only ever reads \code{self}, never writes through it.

\subsection{Splitting the driver into \code{gpio.h} and \code{gpio.c}}
\label{c9:a:3}
So far \secref{c9:a:2}'s \code{gpio_t}, \code{gpio_mode_t}, and every \code{gpio_*()} function
have lived together in one code block. In a real project a driver is split across a header and a
source file, the same split \secref{c1:a:6} introduced for \file{sensor.h}/\file{sensor.c},
include guard and all:
\begin{itemize}
  \item The header (\file{gpio.h}) is the public interface: the \code{gpio_t} and
    \code{gpio_mode_t} type definitions, plus a documented prototype for every function callers
    are allowed to use.
  \item The source file (\file{gpio.c}) \code{#include}s its own header, after the standard and
    AVR headers as everywhere in this book, then defines those same functions. Callers only ever
    see the header; the definitions in \file{gpio.c} are what actually gets compiled.
\end{itemize}
One thing is new compared to \file{sensor.h}: that header only ever declared a function, callers
never touched anything but the value it returned. \file{gpio.h} also declares \code{gpio_t}
itself, because a caller needs to create and initialize their own \code{gpio_t} instances directly
(as in \secref{c9:a:2}'s \code{led1} example), not just receive values back from an opaque handle.
That's also why \code{gpio_t}'s fields stay public here instead of hidden behind a pointer to an
incomplete type; Chapter~\ref{c10:ch} revisits that trade-off.

Here's \file{driver/gpio.h}:

\begin{ccode}
/**
 * @file GPIO driver.
 */
#ifndef DRIVER_GPIO_H_
#define DRIVER_GPIO_H_

#include <stdbool.h>
#include <stdint.h>

/**
 * @brief GPIO driver structure.
 */
typedef struct
{
    /** Pointer to the data direction register associated with the GPIO. */
    volatile uint8_t* ddrx;

    /** Pointer to the output register associated with the GPIO. */
    volatile uint8_t* portx;

    /** Pointer to the input register associated with the GPIO. */
    volatile uint8_t* pinx;

    /** Bit position of the GPIO pin within the associated port. */
    const uint8_t pin;
} gpio_t;

/**
 * @brief Enumeration of GPIO modes.
 */
typedef enum
{
    GPIO_MODE_INPUT,        ///< Standard GPIO input.
    GPIO_MODE_INPUT_PULLUP, ///< GPIO input with its internal pull-up enabled.
    GPIO_MODE_OUTPUT,       ///< GPIO output.
} gpio_mode_t;

/**
 * @brief Initialize GPIO.
 *
 * @param[in] self GPIO instance.
 * @param[in] mode GPIO mode.
 */
void gpio_init(gpio_t* self, gpio_mode_t mode);

/**
 * @brief Set state of the given GPIO.
 *
 * @param[in] self GPIO instance.
 * @param[in] state GPIO state (true = high, false = low).
 *
 * @note This operation is only supported for GPIOs configured as outputs.
 */
void gpio_write(gpio_t* self, bool state);

/**
 * @brief Read state of the given GPIO.
 *
 * @param[in] self GPIO instance.
 *
 * @return GPIO state (true = high, false = low).
 */
bool gpio_read(const gpio_t* self);

/**
 * @brief Toggle state of the given GPIO.
 *
 * @param[in] self GPIO instance.
 */
void gpio_toggle(gpio_t* self);

#endif /** DRIVER_GPIO_H_ */
\end{ccode}

\noindent Here's \file{driver/gpio.c}:

\begin{ccode}
/**
 * @file GPIO driver implementation details.
 */
#include <stdbool.h>
#include <stddef.h>
#include <stdint.h>

#include <avr/io.h>

#include "driver/gpio.h"

// -----------------------------------------------------------------------------
void gpio_init(gpio_t* self, const gpio_mode_t mode)
{
    // Check the GPIO instance, terminate if invalid.
    if (NULL == self) { return; }

    // Configure GPIO mode as specified.
    switch (mode)
    {
        case GPIO_MODE_INPUT_PULLUP:
            // Configure GPIO as input with its internal pull-up enabled.
            *self->ddrx &= ~(1U << self->pin);
            *self->portx |= (1U << self->pin);
            break;
        case GPIO_MODE_OUTPUT:
            // Configure GPIO as output.
            *self->ddrx |= (1U << self->pin);
            *self->portx &= ~(1U << self->pin);
            break;
        default:
            // Configure GPIO as standard input.
            *self->ddrx &= ~(1U << self->pin);
            *self->portx &= ~(1U << self->pin);
            break;
    }
}

// -----------------------------------------------------------------------------
void gpio_write(gpio_t* self, const bool state)
{
    // Check the GPIO instance, terminate if invalid.
    if (NULL == self) { return; }

    // Check GPIO direction, terminate if not output.
    const bool output = (*self->ddrx & (1U << self->pin));
    if (!output) { return; }

    // Set GPIO state as specified.
    if (state) { *self->portx |= (1U << self->pin); }
    else { *self->portx &= ~(1U << self->pin); }
}

// -----------------------------------------------------------------------------
bool gpio_read(const gpio_t* self)
{
    // Return GPIO state if GPIO instance is valid, else false.
    return (NULL != self) ? (*self->pinx & (1U << self->pin)) : false;
}

// -----------------------------------------------------------------------------
void gpio_toggle(gpio_t* self)
{
    // Check the GPIO instance, terminate if invalid.
    if (NULL == self) { return; }

    // Check GPIO direction, terminate if not output.
    const bool output = (*self->ddrx & (1U << self->pin));
    if (!output) { return; }

    // Toggle GPIO state.
    *self->pinx = (1U << self->pin);
}
\end{ccode}

\noindent Notice \code{gpio_init()}'s \code{mode} parameter (and \code{gpio_write()}'s
\code{state}) drop the \code{const} in the header but keep it in the source file. A top-level
\code{const} on a parameter only constrains what the function body can do with its own local copy;
it isn't part of the function's type and can't affect how callers call it, so by convention it's
omitted from the prototype and kept in the definition, where it actually matters.

\subsection{Worked example: an array of LEDs}
\label{c9:a:4}
The whole point of a driver struct is that \code{gpio_write()}/\code{gpio_read()} don't care which
GPIO they're called on. An array of \code{gpio_t} instances plus a loop replaces a whole family of
per-LED macros:

\begin{ccode}
#define LED_COUNT 5U

// Create five LED instances.
static gpio_t leds[LED_COUNT] =
{
    { &DDRD, &PORTD, &PIND, 6U}, // D6  -> PORTD6.
    { &DDRD, &PORTD, &PIND, 7U}, // D7  -> PORTD7.
    { &DDRB, &PORTB, &PINB, 0U}, // D8  -> PORTB0.
    { &DDRB, &PORTB, &PINB, 1U}, // D9  -> PORTB1.
    { &DDRB, &PORTB, &PINB, 2U}, // D10 -> PORTB2.
};

/**
 * @brief Set up system.
 */
static void setup(void)
{
    for (uint8_t i = 0U; i < LED_COUNT; i++)
    {
        gpio_init(&leds[i], GPIO_MODE_OUTPUT);
    }
}
\end{ccode}

\noindent Compare this to \secref{c4:app:b}'s five-LED exercise, Exercise~\ref{c4:ex:7}, which
needed a hand-written \code{LED1_ON} through \code{LED5_ON} (and \code{_OFF}), one pair of macros
per pin. Here, adding a sixth LED on the same port means adding one line to the \code{leds[]}
array, no new functions.

\subsection{Worked example: \code{gpio_t} wired into Chapters 6--8's demo}
\label{c9:a:5}
Combines Chapter~\ref{c6:ch}'s serial driver, Chapter~\ref{c7:ch}'s EEPROM driver, and
Chapter~\ref{c8:ch}'s watchdog driver with this section's \code{gpio_t}: LED1 (D9) persists its
state across resets via EEPROM, BTN1 (D13) toggles LED1 through a debounced pin-change interrupt
exactly like \secref{c8:a:5}, and the watchdog resets the chip if the main loop ever stalls. Every
register access that went through Chapter~\ref{c8:ch}'s \code{LED1_TOGGLE}/\code{LED1_ENABLED}/%
\code{BTN1_PRESSED} macros now goes through \code{gpio_toggle()}/\code{gpio_read()} instead.

The circuit is Chapter~\ref{c6:ch}'s, in Figure~\ref{c6:fig:serial}.

\begin{console}
gpio_driver_demo/
├── Makefile
├── main.c
├── include/
│   └── driver/
│       ├── eeprom.h
│       ├── gpio.h
│       ├── serial.h
│       └── watchdog.h
└── source/
    └── driver/
        ├── eeprom.c
        ├── gpio.c
        ├── serial.c
        └── watchdog.c
\end{console}

\begin{itemize}
  \item The serial, EEPROM, and watchdog drivers were introduced in
    Chapters~\ref{c6:ch}--\ref{c8:ch}.
  \item \file{driver/gpio.h} and \file{driver/gpio.c}, both printed in full below, extend the
    \code{gpio_t} type introduced in \secref{c9:a:2} and \secref{c9:a:3}.
  \item The extended driver adds one capability not covered in those sections:
  \begin{itemize}
    \item \code{gpio_enable_pci()} enables or disables the pin-change interrupt for a GPIO.
    \item It does this by writing the same \code{PCICR} and \code{PCMSKx} bits that were
      configured manually in \secref{c8:a:5}.
  \end{itemize}
  \item Supporting \code{gpio_enable_pci()} requires two additional fields in \code{gpio_t},
    beyond \code{ddrx}, \code{portx}, \code{pinx}, and \code{pin}:
  \begin{itemize}
    \item A pointer to the port's \code{PCMSKx} register.
    \item The bit position of the port's \code{PCIEx} bit in \code{PCICR}.
  \end{itemize}
  \item Every \code{gpio_t} instance contains these two additional fields:
  \begin{itemize}
    \item This is true even when that particular GPIO never uses \code{gpio_enable_pci()}.
    \item For example, LED1 still carries the pin-change-interrupt configuration fields.
  \end{itemize}
\end{itemize}

\subsubsection{The driver in full}
\cfile{lectures/L09/gpio_driver_demo/include/driver/gpio.h}
\needspace{8\baselineskip}
\cfile{lectures/L09/gpio_driver_demo/source/driver/gpio.c}

\subsubsection{\code{main.c} in full}
\cfile{lectures/L09/gpio_driver_demo/main.c}

\noindent\code{led1_event}, the debounce timer, and the EEPROM-persistence logic are unchanged from
\secref{c8:a:5}; only the GPIO handling changed, from bespoke macros tied to \code{PORTB}
specifically to a driver that works for however many \code{gpio_t} instances a program declares,
on whichever pins and ports they happen to be wired to.

\subsubsection{Building and running}
Build and flash the same way as \secref{c8:a:5}: on Linux, run \code{make} in the demo's
directory, \file{lectures/L09/gpio_driver_demo} (\code{make flash} to program the board); on
Windows via Microchip Studio, recreate this folder structure inside the project and add
\file{include} as an include directory.

The demo's \file{Makefile} is the one printed in \secref{c6:sec:build}, with
\file{source/driver/eeprom.c}, \file{source/driver/gpio.c} and \file{source/driver/watchdog.c}
added to its \code{SRC_FILES}. Its serial, EEPROM and watchdog drivers are byte for byte the ones
printed in \secref{c6:a:6}, \secref{c7:a:4} and \secref{c8:a:5}.

% Section 9.2, from lectures/L09/README.md: the goals, the evaluation questions, and the next
% lecture.

\needspace{10\baselineskip}
\section{Review}
\label{c9:sec:review}

\subsection*{What you should be able to do now}
After this chapter, you should:
\begin{itemize}
  \item Explain why hardcoded per-instance functions or macros break down once a program needs
    more than one of something.
  \item Be able to design a struct that bundles a peripheral's configuration and state.
  \item Be able to write ``methods'' as plain functions taking a pointer to that struct as their
    first parameter (the \code{self} convention).
  \item Refactor a program built on duplicated globals into one built on an array of struct
    instances.
\end{itemize}

\needspace{6\baselineskip}
\subsection*{Questions to test yourself}
\begin{enumerate}
  \item What problem does bundling a pin number into a struct solve that a \code{#define}d macro
    can't?
  \item Why does \code{gpio_read()} take a pointer to a \code{gpio_t} rather than a \code{gpio_t}
    by value?
  \item How would you rewrite five separate
    \code{LED1_ON}/\allowbreak\code{LED2_ON}/\allowbreak\dots{} macros as a single
    \code{gpio_write()} function plus an array of \code{gpio_t} instances?
\end{enumerate}

\needspace{6\baselineskip}
\subsection*{Looking ahead}
Chapter~\ref{c10:ch}: hiding a struct's implementation details (opaque structs) and manual
polymorphism via function pointers and vtables.

% Section 9.3, from lectures/L09/appendix/b_exercises.md.

\section{Exercises}
\label{c9:sec:exercises}\label{c9:app:b}

\begin{aside}
\textbf{How to work these.} The Reflection exercises are answered on paper, from the chapter. The
Code exercises are programs for the board: build and flash them the way \secref{c2:app:c}
(Microchip Studio) or \secref{c2:app:d} (Linux) describes, or run them in the simulator of
Appendix~\ref{sim:ch} if you have no board yet.

\textbf{How to check your work.} The Reflection exercises, and the questions Code
Exercises~\ref{c9:ex:3} to~\ref{c9:ex:6} ask about their design, are answered in
Appendix~\ref{app:answers}. The Code exercises have no printed solution: the board tells you
whether a program does what the exercise describes. Try each exercise before you read its answer.
\end{aside}

% ------------------------------------------------------------------------------------------
\exerciseset{Concepts}
\label{c9:set:1}

\reflexercise{A struct or a macro}
\label{c9:ex:1}
What problem does bundling a pin number into a \code{gpio_t} struct solve that a
\code{#define}d \code{LED_ON} macro can't? What has to change to support a second LED under each
approach?

\reflexercise{By pointer or by value}
\label{c9:ex:2}
\code{gpio_read()} takes \code{const gpio_t* self} rather than \code{gpio_t self} (the whole
struct, by value). What would change, functionally or in cost, if it took \code{gpio_t self}
instead? Would \code{gpio_init()} still work correctly if it took \code{self} by value?

% ------------------------------------------------------------------------------------------
\exerciseset{Program}
\label{c9:set:2}

\codeexercise{Freeing a GPIO}
\label{c9:ex:3}
Add \code{gpio_free(gpio_t* self)} to \secref{c9:a:2}'s driver:
\begin{itemize}
  \item Reset the pin's \code{DDRx}/\code{PORTx} bits back to their power-on state (input,
    pull-up disabled).
  \item Decide whether \code{self->ddrx}, \code{self->portx}, and \code{self->pinx} should be set
    to \code{NULL} afterwards.
  \item Justify the choice:
  \begin{itemize}
    \item What does calling \code{gpio_write()} or \code{gpio_toggle()} on a freed instance do in
      each case?
    \item Which is preferable?
  \end{itemize}
\end{itemize}

\codeexercise{Initializing from a pin number}
\label{c9:ex:4}
Rewrite \code{gpio_init()} so callers pass a single pin number, 0--13, instead of
\code{DDRx}/\code{PORTx}/\code{PINx} pointers:
\begin{itemize}
  \item Pin mapping:
  \begin{itemize}
    \item 0--7 map to \code{PORTD0}--\code{PORTD7}.
    \item 8--13 map to \code{PORTB0}--\code{PORTB5}.
  \end{itemize}
  \item \code{gpio_init()} should compute and store \code{ddrx}/\code{portx}/\code{pinx} (and the
    bit position) itself.
  \item Return \code{bool}:
  \begin{itemize}
    \item \code{false} for any pin number outside 0--13, without touching any register.
    \item \code{true} otherwise, after initialization.
  \end{itemize}
  \item \code{pin} is \code{const} in \secref{c9:a:2}'s \code{gpio_t} because it's only ever set
    at struct-literal time; now that \code{gpio_init()} derives it from the pin number instead,
    does it still make sense for \code{pin} to stay \code{const}?
\end{itemize}

\codeexercise{A pin-reservation table}
\label{c9:ex:5}
Add a pin-reservation table:
\begin{itemize}
  \item Declare a \code{static uint16_t pin_reg} in \file{gpio.c}.
  \item Use one bit for each pin number introduced in Exercise~\ref{c9:ex:4}:
  \begin{itemize}
    \item Bit~$n$ represents pin~$n$.
    \item \code{1} means that the pin is reserved.
    \item \code{0} means that the pin is free.
  \end{itemize}
  \item Update \code{gpio_init()}:
  \begin{itemize}
    \item Check the bit corresponding to the requested pin.
    \item If the bit is already set, return \code{false}.
    \item In that case, do not modify any hardware register.
    \item If the bit is clear, set it to reserve the pin.
    \item Then initialize the GPIO as in Exercise~\ref{c9:ex:4}.
    \item Return \code{true} when the initialization succeeds.
  \end{itemize}
  \item Extend the \code{gpio_free()} function from Exercise~\ref{c9:ex:3}:
  \begin{itemize}
    \item Clear the pin's bit in \code{pin_reg}.
    \item This allows a later call to \code{gpio_init()} to reserve the same pin again.
  \end{itemize}
  \item Answer the following questions:
  \begin{itemize}
    \item Why must \code{pin_reg} be declared \code{static}?
    \item What would break if it were not \code{static}?
  \end{itemize}
\end{itemize}

\codeexercise{A timer driver}
\label{c9:ex:6}
Capstone --- harder: Implement a timer driver named \code{timer_t}:
\begin{itemize}
  \item Follow the header/source separation used in \secref{c9:a:3}:
  \begin{itemize}
    \item \file{driver/timer.h}
    \item \file{driver/timer.c}
  \end{itemize}
  \item The timer must run a callback function whenever its timeout expires.
  \item \code{timer_t} should follow the same ownership model as \code{gpio_t}:
  \begin{itemize}
    \item The caller declares and owns each \code{timer_t} instance.
    \item The caller may declare an array of timer instances, just as \secref{c9:a:4} declares an
      array of \code{gpio_t} instances.
    \item \code{timer_init()} initializes an existing instance in place.
  \end{itemize}
  \item The ATmega328P has only three hardware timer/counters:
  \begin{itemize}
    \item At most three \code{timer_t} instances may therefore be active at the same time.
    \item Use a reservation scheme similar to Exercise~\ref{c9:ex:5}'s \code{pin_reg}.
    \item \code{timer_init()} must not create more active timer instances than there are hardware
      timer/counters available.
  \end{itemize}
  \item \code{timer_t} is a complete timer driver instance:
  \begin{itemize}
    \item Each instance has:
    \begin{itemize}
      \item Its own callback.
      \item Its own elapsed-time count.
      \item Its own start/stop state.
      \item One reserved hardware timer/counter.
    \end{itemize}
    \item It is driven directly by the hardware timer/counter and its compare-match interrupt.
  \end{itemize}
\end{itemize}

\task{Recommended \file{driver/timer.h}}
\code{timer_callback_t} is a function pointer, see \secref{c1:a:8} if that syntax is unfamiliar.
Document every declaration with Doxygen comments, the same way
\secref{c9:a:2}/\secref{c9:a:3}'s \file{gpio.h} does, and add the header guard yourself (suggested
guard: \code{DRIVER_TIMER_H_}).

\begin{ccode}
#include <stdbool.h>
#include <stdint.h>

/** Alias for timer callbacks. */
typedef void (*timer_callback_t)(void);

typedef struct
{
    // Left for you to design.
} timer_t;

bool timer_init(timer_t* self, uint32_t timeout_ms, timer_callback_t callback);
void timer_clear(timer_t* self);
bool timer_running(const timer_t* self);
void timer_start(timer_t* self);
void timer_stop(timer_t* self);
void timer_toggle(timer_t* self);
void timer_reset(timer_t* self);
\end{ccode}

\task{Initialization and reservation}
\begin{itemize}
  \item \code{timer_init()} must return \code{false} when:
  \begin{itemize}
    \item \code{timeout_ms == 0}.
    \item \code{callback == NULL}.
    \item No hardware timer/counter is available.
  \end{itemize}
  \item Otherwise, \code{timer_init()} must:
  \begin{itemize}
    \item Reserve one free hardware timer/counter.
    \item Associate it with \code{self}.
    \item Store the timeout and callback.
    \item Initialize the timer instance.
    \item Return \code{true}.
  \end{itemize}
  \item \code{timer_clear()} must:
  \begin{itemize}
    \item Stop the timer.
    \item Release its reserved hardware timer/counter.
    \item Allow that hardware timer/counter to be claimed by a later \code{timer_init()} call.
    \item Work whether the later call uses the same \code{timer_t} instance or another one.
  \end{itemize}
  \item Consider the same design question raised for \code{gpio_free()} in
    Exercise~\ref{c9:ex:3}:
  \begin{itemize}
    \item Should \code{timer_clear()} also reset the fields of \code{self}?
    \item If so, which values would leave the instance in a safe, inert state?
  \end{itemize}
\end{itemize}

\task{Timer control}
\begin{itemize}
  \item \code{timer_start()} must start or enable the timer.
  \begin{itemize}
    \item It must not change the elapsed-time count.
  \end{itemize}
  \item \code{timer_stop()} must stop or disable the timer.
  \begin{itemize}
    \item It must not change the elapsed-time count.
  \end{itemize}
  \item \code{timer_toggle()} must switch between the running and stopped states.
  \begin{itemize}
    \item It must not change the elapsed-time count.
  \end{itemize}
  \item \code{timer_reset()} must:
  \begin{itemize}
    \item Set the elapsed-time count to zero.
    \item Preserve the timer's current running or stopped state.
  \end{itemize}
\end{itemize}

\task{Driver independence}
\begin{itemize}
  \item The timer driver must not know that \code{gpio_t} exists.
  \item \file{timer.c} must not directly control an LED or read a button.
  \item Connecting a timer callback to an LED is the responsibility of \file{main.c}.
\end{itemize}

\task{Hardware demonstration}
Wire the following circuit:
\begin{itemize}
  \item Connect a button to D13.
  \item Connect an LED to D8.
  \item Configure a timer with a timeout of 100\,ms.
\end{itemize}
The program must behave as follows:
\begin{itemize}
  \item Pressing the button toggles the timer between running and stopped.
  \item While the timer is running:
  \begin{itemize}
    \item Its callback runs every 100\,ms.
    \item The callback toggles the LED.
  \end{itemize}
  \item While the timer is stopped:
  \begin{itemize}
    \item The LED remains off.
    \item Calling the callback manually must have no effect.
  \end{itemize}
  \item To enforce this behavior, the callback itself must check \code{timer_running()} before
    toggling the LED.
\end{itemize}

\task{Testing}
Confirm \code{timer_init()} returns \code{false} for each of the three invalid cases in
``Initialization and reservation'': a \code{0} timeout, a \code{NULL} callback, and a fourth
\code{timer_t} once three hardware timer/counters are already reserved. Then confirm
\code{timer_clear()} actually frees a hardware timer/counter for reuse: after reserving all three,
clear one and check that a new \code{timer_init()} call succeeds again.

